\section*{第 \ref{ch:b3:submanifolds} 章——
$\R^n$ 中的子流形}

\begin{solution}{exo:b3:submanifolds:1}
(a) 各集合都是某淹没的 $F^{-1}(0)$：取 $\norm x^2 - 1$ 在
$\R^n\setminus\{0\}$ 上（梯度 $2x \neq 0$）：维数 $n-1$；
$x^2 + y^2 - z^2 - 1$（梯度 $(2x, 2y, -2z) \neq 0$ 于
单叶双曲面上，其中 $x^2 + y^2 = 1 + z^2 > 0$）：维数 $2$；
环面函数 $G = (\rho - R)^2 + z^2 - r^2$，$\rho =
\sqrt{x^2+y^2}$，在环面附近为 $\mathcal C^\infty$（$\rho
\geq R - r > 0$ 处成立），且 $\nabla G \neq 0$（其
$z$ 分量为 $2z$；在 $z = 0$ 处，径向分量
为 $2(\rho - R)\ne0$，因 $\abs{\rho - R} = r$）：维数
$2$。

(b) 当 $\varepsilon$ 充分小时，
$(C\setminus\{0\})\cap B(0,\varepsilon)$ 恰有 $2$
个\omterm{def:b3:topology:components}{连通分支}（上、下两叶去心锥面，每叶
\omterm{def:b3:topology:pathconnected}{道路}连通：经由圆周与射线相连）。若 $C$
为 $2$ 维子流形（在 $0$ 处），则拉直可得一个
从 $C\cap\Omega$ 到平面某开片上的\omterm{def:b3:topology:continuity}{同胚}，
把 $0$ 映为点 $p$；
$p$ 的小去心平面\omterm{def:b3:topology:topology}{邻域}只有\emph{一个}\omterm{def:b3:topology:components}{连通分支}，而
\omterm{def:b3:topology:continuity}{同胚}保持去心
\omterm{def:b3:topology:topology}{邻域}的\omterm{def:b3:topology:components}{连通分支}个数，矛盾。
\end{solution}

\begin{solution}{exo:b3:submanifolds:2}
(a) $T_aS^{n-1} = \ker\bigl(2a^{\mathsf T}\bigr) = a^\perp$。
(b) 在 $p = ((R+r)\cos\theta, (R+r)\sin\theta, 0)$ 处：$\nabla G
= (2r\cos\theta, 2r\sin\theta, 0)$，于是切平面为
$\operatorname{Vect}\bigl((-\sin\theta, \cos\theta, 0),\ (0,
0, 1)\bigr)$：与外侧赤道相切的
竖直平面。
(c) 螺旋线是嵌入曲线，且 $\varphi'(t_0) =
(-\sin t_0, \cos t_0, 1) \neq 0$，
$\varphi(t_0)$ 处的切线为 $\varphi(t_0) + \R\,\varphi'(t_0)$，如
\cref{prop:b3:submanifolds:tangent}(2) 所给出。
\end{solution}

\begin{solution}{exo:b3:submanifolds:3}
(a) $Df(x,y) = \bigl(\begin{smallmatrix}2x & -2y\\ 2y &
2x\end{smallmatrix}\bigr)$，$\det = 4(x^2 + y^2)$：
除原点外，逆函数定理在每一点处都适用。
(b) $f(-z) = f(z)$：在任何关于
$0$ 对称的集合上都不是单射；在 $\R^2\setminus\{0\}$ 上处处为局部微分\omterm{def:b3:topology:continuity}{同胚}，
但整体上是 $2$ 对 $1$ 的。在 $(1, 0) = f(\pm(1,
0))$ 附近，两个逆映射即两支平方根，
用复数记号写为 $w \mapsto \pm\sqrt w$（主支）：
即
\[
(u, v) \longmapsto \pm\Bigl(\sqrt{\tfrac{u +
\sqrt{u^2+v^2}}{2}},\ \frac{v}{2\sqrt{(u +
\sqrt{u^2+v^2})/2}}\Bigr).
\]
(c) 雅可比行列式为 $r > 0$：在
$\intoo0\infty\times\R$ 上为局部微分\omterm{def:b3:topology:continuity}{同胚}，但 $\theta \mapsto \theta + 2\pi$
给出同一点：局部可逆（角度在
半平面上只确定到 $2\pi$ 之差），整体上绝不可逆。
\end{solution}

\begin{solution}{exo:b3:submanifolds:4}
(a) $\nabla F = 3(x^2 - y,\ y^2 - x)$ 为零当且仅当 $y = x^2$
且 $x = y^2$，即 $x^4 = x$：驻点为 $(0,0)$ 与 $(1,1)$；仅
$(0,0)$ 在 $\mathcal C$ 上（$F(1,1) = -1$）。于是在 $\mathcal
C\setminus\{0\}$ 上 $F$ 为淹没，故为 $1$ 维子流形，
局部为 $y = \psi(x)$ 的图像，其中 $F_y = 3(y^2 - x) \neq 0$，
或 $x = \chi(y)$，其中 $F_x = 3(x^2 - y) \neq 0$（
原点之外二者至少成立其一）。

(b) 在 $(\frac32, \frac32)$ 处：$\nabla F = 3(\frac94 - \frac32)
(1, 1) = \frac94(1,1)$，切线为 $x + y = 3$。

(c) 有理参数化 $x = \frac{3t}{1 + t^3}$，
$y = \frac{3t^2}{1+t^3}$ 过 $0$，且在 $t = 0$ 处速度为
$(3, 0)$；交换 $x \leftrightarrow y$（
曲线对称，或取 $1/t$ 作重新参数化）得到第二条
$\mathcal C^1$ 曲线过 $0$，速度为 $(0, 3)$。两个
独立的切方向对
$1$ 维子流形是不可能的（其\omterm{def:b3:submanifolds:tangent}{切空间}为直线，
见\cref{prop:b3:submanifolds:tangent}）：$\mathcal C$ 在原点处不是
子流形——一个横截的自交点。
\end{solution}

\begin{solution}{exo:b3:submanifolds:5}
(a) $F(M + H) = M^{\mathsf T}M + M^{\mathsf T}H + H^{\mathsf
T}M + H^{\mathsf T}H$：$DF(M)(H) = M^{\mathsf T}H +
H^{\mathsf T}M$。设 $M \in O_n$，$S$ 为对称矩阵，$H =
\frac12MS$，则 $DF(M)(H) = \frac12(S + S^{\mathsf T}) = S$；
满射到 $S_n$ 上。

(b) $O_n = F^{-1}(I)$，$F$ 为淹没（到 $S_n$ 上，
维数 $\frac{n(n+1)}2$）在其每点处：
为子流形，维数为 $n^2 - \frac{n(n+1)}2 =
\frac{n(n-1)}2$。紧：闭（$F$ 连续），且
有界（各列为单位向量）。$I$ 处的\omterm{def:b3:submanifolds:tangent}{切空间}为 $\ker DF(I) =
\{H : H + H^{\mathsf T} = 0\}$。

(c) $\bigl(\eu^{tH}\bigr)^{\mathsf T}\eu^{tH} =
\eu^{tH^{\mathsf T}}\eu^{tH} = \eu^{-tH}\eu^{tH} = I$
（转置可穿过级数；可交换矩阵的指数
满足乘法法则，
见\cref{thm:b3:ode:matrixexp}）。
\end{solution}

\begin{solution}{exo:b3:submanifolds:6}
(a) $\det(M + H) = \det M\,\det(I + M^{-1}H) = \det M\bigl(1
+ \operatorname{tr}(M^{-1}H) + O(\norm H^2)\bigr)$，对
可逆矩阵 $M$（$\det$ 在 $I$ 附近展开式中的线性项
$\prod(1 + \lambda_i)$）；记 $\det M\cdot M^{-1} =
\operatorname{com}(M)^{\mathsf T}$：
$D\det(M)(H) = \operatorname{tr}\bigl(\operatorname{com}
(M)^{\mathsf T}H\bigr)$，且该公式对所有 $M$
亦成立（由稠密性与连续性）。在 $SL_n$ 上，$\det M = 1$：
$D\det(M) \ne 0$（它在 $H = M$ 处的值为
$\operatorname{tr}(I)\cdot\dots = n\det M = n$）。

(b) $SL_n = \det^{-1}(1)$，其中 $\det$ 在该处为淹没
（取值于 $\R$）：维数为 $n^2 - 1$；$T_ISL_n = \ker
D\det(I) = \{H : \operatorname{tr}H = 0\}$。

(c) $GL_n$ 是 $M_n(\R)$ 的\emph{开}子集
（\cref{exo:b3:topology:8}）：是满维的子流形，
维数为 $n^2$（拉直：取恒等坐标卡）。
\end{solution}

\begin{solution}{exo:b3:submanifolds:7}
(a) $(y, x) = \lambda(2x, 2y)$ 且 $x^2 + y^2 = 1$：$y =
2\lambda x$，$x = 2\lambda y$ 可得 $x^2 = y^2 = \frac12$。
$xy$ 的取值为 $\pm\frac12$：最大值 $\frac12$，在
$\pm\frac1{\sqrt2}(1,1)$ 处取得；最小值 $-\frac12$，在
$\pm\frac1{\sqrt2}(1,-1)$ 处取得（约束集紧：
极值存在）。

(b) 在单纯形\omterm{def:b3:topology:interior}{内部}（$p_i > 0$）上，用拉格朗日乘数法求
$H(p) = -\sum p_i\ln p_i$ 在约束 $\sum p_i = 1$ 下的极值：
$-\ln p_i - 1 = \lambda$（对所有 $i$），所有 $p_i$ 相等，$p_i =
\frac1n$，此时 $H = \ln n$。在紧的
单纯形上最大值存在。若其在边界上取得，
（某个 $p_i = 0$），则分布仅位于 $\leq n - 1$ 个点上，
由归纳法 $H \leq \ln(n-1) < \ln n$，
故\omterm{def:b3:topology:interior}{内部}临界点为全局最大值——
均匀的无知使熵最大。

(c) 在紧的椭圆上极小化 $(x-1)^2 + y^2$：
$(2(x{-}1), 2y) = \lambda(\frac x2, 2y)$。若 $y \neq 0$，
则 $\lambda = 1$，由 $2(x - 1) = \frac x2$ 得 $x =
\frac43$，$y^2 = 1 - \frac49 = \frac59$，距离$^2$ $=
\frac19 + \frac59 = \frac23$。若 $y = 0$，则 $x = \pm2$，
距离分别为 $1$ 与 $3$。最近点为 $\bigl(\frac43,
\pm\frac{\sqrt5}3\bigr)$，距离为 $\sqrt{2/3} < 1$。
乘子方程表明，从 $(1,0)$ 到最近点的线段
与 $\nabla$（椭圆）平行：它与椭圆
正交相交，正如几何所要求的。
\end{solution}

\begin{solution}{exo:b3:submanifolds:8}
对 $n$ 作归纳；$n = 1$ 时显然。瑞利函数
$f(x) = \langle Ax, x\rangle$ 取得其最大值 $\lambda_1$
于紧集 $S^{n-1}$ 上的某点 $v_1$ 处；由拉格朗日乘数法
（\cref{thm:b3:submanifolds:lagrange}，球面作为水平集）
可得 $2Av_1 = 2\lambda v_1$，且 $\lambda = \langle Av_1,
v_1\rangle = \lambda_1$。超平面 $v_1^\perp$ 关于
$A$ 不变（由对称性，$\langle Av, v_1\rangle = \langle
v, Av_1\rangle = 0$）；限制后仍为对称算子，且
由归纳假设得标准正交特征基 $v_2, \dots, v_n$，
位于 $v_1^\perp$ 中，其特征值为 $\lambda_2 \geq \dots \geq
\lambda_n$，每个均为 $f$ 在
其余正交补的球面上的最大值。库朗–费歇尔定理的证明
与\cref{exo:b3:spectral:8}中完全相同：将 $x = \sum c_iv_i$ 展开；
在 $k$ 维检验空间上，与
$\operatorname{Vect}(v_k, \dots, v_n)$ 相交（维数计数
在 $\R^n$ 中），得 $\min \leq \lambda_k$，而
$\operatorname{Vect}(v_1, \dots, v_k)$ 取得 $\min =
\lambda_k$。
\end{solution}

\begin{solution}{exo:b3:submanifolds:9}
将每列规范化为单位范数，使 $\abs{\det}$ 除以
$\prod\norm{c_i}$，只需证明 $\abs{\det M} \leq 1$，
在各列均为单位向量时成立，且等号当且仅当 $M \in O_n$。
函数 $\det$ 在紧集
$(S^{n-1})^n$ 上连续，取得最大值 $m \geq \det I = 1 > 0$，
于某点 $M$ 处取得。固定除第 $i$ 列外的所有列，$\det$ 关于
$c_i$ 为线性，梯度为第 $i$ 列
（取自 $\operatorname{com}(M)$）；用拉格朗日乘数法于第 $i$ 个球面上：
$\operatorname{com}(M)_{\cdot i} = \lambda_i c_i$。
由恒等式 $M^{\mathsf T}\operatorname{com}(M) = \det(M)\,I$
可写为 $\langle c_j, \operatorname{com}(M)_{\cdot i}\rangle =
\det(M)\,\delta_{ij}$，即 $\lambda_i\langle c_j,
c_i\rangle = \det(M)\delta_{ij}$；取 $j = i$，得 $\lambda_i
= \det M = m \neq 0$，再取 $j \neq i$ 得 $\langle c_i,
c_j\rangle = 0$，即各列标准正交，$M \in O_n$，
$m = \abs{\det M} = 1$。因此 $\abs{\det} \leq
\prod\norm{c_i}$ 恒成立，且等号恰当列正交时取得
（再反向缩放）：平行多面体的体积在
给定棱长时，于各棱两两垂直时最大。
\end{solution}

\begin{solution}{exo:b3:submanifolds:10}
$y = \pm\sqrt{1 - x^2}$ 在每个 $y \neq
0$ 的点附近成立；$x = \pm\sqrt{1 - y^2}$ 在每个 $x \neq 0$ 的点附近成立；
在 $(1, 0)$ 处：图像 $x = \sqrt{1 - y^2}$ 在 $y \in
\intoo{-1}1$ 上给出，它即为
\cref{thm:b3:submanifolds:characterizations}(3) 的
坐标互换情形。总有某种排列可行，因为
一维的切线不可能同时竖直又水平——但
可以是二者之一，故
对“因变量”坐标的固定选择
并不能在每一点都适用。
\end{solution}

\begin{solution}{exo:b3:submanifolds:11}
(a) 定义映射 $F(M) = M^{\mathsf T}M - I$ 连续：
$O_n = F^{-1}(0)$ 为闭集；正交矩阵的各列
为单位向量，故其欧几里得
（弗罗贝尼乌斯）范数恰为 $\sqrt n$：有界。
由海涅–博雷尔定理，在 $M_n(\R) \cong \R^{n^2}$ 中有界闭集为紧。

(b) $O_n$ 是本章所研究的水平集 $F = 0$：
$DF(A)H = A^{\mathsf T}H + H^{\mathsf T}A$，
在每个 $A \in O_n$ 处满射至对称矩阵（对任意对称 $S$，
取 $H = \frac12AS$），故 $O_n$ 是
维数为 $n^2 - \frac{n(n+1)}2 = \frac{n(n-1)}2$，且
\[
T_AO_n = \ker DF(A) = \{H : A^{\mathsf T}H \text{
反对称}\} = \{AK : K^{\mathsf T} = -K\} ;
\]
当 $A = I$ 时，它们即为反对称矩阵。

(c) $\exp(tK)^{\mathsf T}\exp(tK) = \exp(tK^{\mathsf T})
\exp(tK) = \exp(-tK)\exp(tK) = I$（两矩阵 $\pm tK$
可交换，故指数之积等于
和的指数）：$\exp(tK) \in O_n$，且 $\gamma(t) =
A\exp(tK)$ 为 $O_n$ 中的曲线，满足 $\gamma(0) = A$，
$\gamma'(0) = AK$。当 $K$ 遍历反对称矩阵时，
$AK$ 扫过 $T_AO_n$：指数映射
通过显式曲线实现整个\omterm{def:b3:submanifolds:tangent}{切空间}——这即李群捷径，
正是 \cref{pb:b3:submanifolds:1} 对 $SO(3)$ 所利用的。
\end{solution}

\begin{solution}{exo:b3:submanifolds:12}
存在性：将 $M$ 与以
$p$ 为中心的大闭球相交，得一非空紧集；连续距离
在其上取得最小值，球外的点都
更远。一阶条件：设 $\gamma$ 为 $M$ 中的曲线，
满足 $\gamma(0) = x_0$，函数 $h(t) = \norm{\gamma(t)
- p}^2$ 可导，且在 $0$ 处取得最小值：
\[
0 = h'(0) = 2\,\langle\gamma'(0),\ x_0 - p\rangle,
\]
而 $\gamma'(0)$ 遍历 $T_{x_0}M$，故 $p - x_0 \perp
T_{x_0}M$。球面 $S(c, r)$ 在 $x_0$ 处的\omterm{def:b3:submanifolds:tangent}{切空间}为
$(x_0 - c)^\perp$，故 $p - x_0 \parallel x_0 - c$：$x_0$
位于过 $c$ 与 $p$ 的直线上，距离为 $r$，
即自 $c$ 出发的射线上的点——几何直观正是如此。抛物线：在
$x_0 = (x, x^2)$ 处的\omterm{def:b3:submanifolds:tangent}{切空间}由 $(1, 2x)$ 张成；
与 $p - x_0 = (2 - x, -x^2)$ 正交的条件为
\[
(2 - x) - 2x^3 = 0,
\qquad\text{即}\qquad 2x^3 + x - 2 = 0,
\]
其唯一实根（$x \mapsto 2x^3 + x$ 严格
递增），故 $x \approx 0.835$；于是 $x_0 \approx (0.835,
0.698)$，且 $d(p, M) = \sqrt{(2 - 0.835)^2 + 0.698^2}
\approx 1.358$。
\end{solution}

\begin{solution}{pb:b3:submanifolds:1}
\textbf{1.} 映射 $q = t + x\mathrm i + y\mathrm j + z\mathrm k
\mapsto \bigl(\begin{smallmatrix}\alpha & \beta\\ -\bar\beta
& \bar\alpha\end{smallmatrix}\bigr)$，其中 $\alpha = t + \iu
x$，$\beta = y + \iu z$：易验证 $1, \mathrm i,
\mathrm j, \mathrm k$ 分别对应 $I$，$\bigl(\begin{smallmatrix}
\iu & 0\\ 0 & -\iu\end{smallmatrix}\bigr)$,
$\bigl(\begin{smallmatrix}0 & 1\\ -1 &
0\end{smallmatrix}\bigr)$, $\bigl(\begin{smallmatrix}0 &
\iu\\ \iu & 0\end{smallmatrix}\bigr)$，其积
复现四元数乘法表：该映射为单射的
代数同态，故 $\mathbb H$ 继承结合律；
$N(q) = \abs\alpha^2 + \abs\beta^2 = \det$ 可乘，
共轭对应于伴随矩阵的转置，
由此得 $\overline{pq} = \bar q\bar p$。中心：
与 $\mathrm i$ 可交换迫使 $y = z = 0$，与 $\mathrm j$ 可交换迫使
$x = 0$，故中心为 $\R$。

\textbf{2.} $q\bar q = N(q)$：当 $q \neq 0$ 时，$q^{-1} = \bar
q/N(q)$，故为除代数。在 $S^3$ 上：$N(pq) = 1$ 且
$N(q^{-1}) = 1$，故为群；且 $S^3 \subseteq \R^4$ 是
单位球面，是紧的 $3$ 维子流形。

\textbf{3.} $v$ 为纯四元数，当且仅当 $\bar v = -v$；此时
$\overline{qv\bar q} = q\bar v\bar q = -qv\bar q$：$\rho_q$
保持 $P$。线性性显然；$N(qv\bar q) =
N(q)N(v)N(q) = N(v)$：是 $(P, N) \cong (\R^3,
\norm\cdot^2)$：$\rho_q \in O(3)$。又 $\rho_{pq}(v) =
pqv\overline{pq} = p(qv\bar q)\bar p =
\rho_p(\rho_q(v))$：同态。

\textbf{4.} $\rho_q = \mathrm{id}$ 当且仅当 $qv = vq$ 对所有纯
$v$，即 $q$ 与 $\mathrm i, \mathrm j, \mathrm k$ 可交换，
亦即 $q$ 位于中心（第 1 问）：$q \in \R\cap S^3 =
\{\pm1\}$。

\textbf{5.} 记 $q = t + p$（$t \in \R$，$p$ 为纯四元数）：$1 =
N(q) = t^2 + N(p)$，故 $t = \cos\frac\theta2$，且 $p =
\sin\frac\theta2\,u$，其中 $N(u) = 1$，对某个 $\theta$
（若 $p = 0$，则 $q = \pm1$ 作用平凡）。因为 $u^2 = -N(u)
= -1$，$q$ 与 $u$ 可交换，且 $\rho_q(u) = qu\bar q =
uq\bar q = u$：此即旋转轴。对满足 $w \perp u$ 的纯单位四元数，
乘法规则 $vw = -\langle v, w\rangle + v\wedge w$
（由乘法表按分量展开）得 $uw = u\wedge w
= -wu$。于是
\[
\rho_q(w) = \bigl(\cos\tfrac\theta2 +
\sin\tfrac\theta2u\bigr)\,w\,\bigl(\cos\tfrac\theta2 -
\sin\tfrac\theta2u\bigr)
= \cos\theta\,w + \sin\theta\,(u\wedge w),
\]
利用 $uwu = -u^2w = w$ 及二倍角公式，即
转角为 $\theta$ 的旋转，位于定向平面 $(w,
u\wedge w)$。

\textbf{6.} 每个 $\rho_q$ 是绕 $u$ 的旋转，转角为
$\theta$：在标准正交基 $(u, w, u\wedge w)$ 下，其
矩阵的行列式为 $+1$：$\operatorname{im}\rho \subseteq
SO(3)$。满射性：对矩阵 $R \in SO(3)$，它以 $1$ 为
特征值，因为
\[
\det(R - I) = \det R\,\det(I - R^{\mathsf T}) = \det(I - R)
= (-1)^3\det(R - I),
\]
故 $\det(R - I) = 0$。取单位特征向量 $u$；$R$
保持 $u^\perp$，并在其上限制为
某转角 $\theta$ 的旋转（平面正交变换，行列式为 $1$）：
$R = \rho_q$，取 $q = \cos\frac\theta2 +
\sin\frac\theta2\,u$。由第 4 问及第一
同构（\cref{thm:b3:groups:firstiso}）：$SO(3)
\cong S^3/\{\pm1\}$。

\textbf{7.} $O_3$ 是紧的 $3$ 维子流形
（见\cref{exo:b3:submanifolds:5}）；$\det$ 在其上连续
取值于 $\{\pm1\}$，故 $SO(3) = O_3\cap\{\det = 1\}$
在 $O_3$ 中既开又闭，是若干\omterm{def:b3:topology:components}{连通分支}的并，
从而它本身是紧的 $3$ 维子流形，
且在 $I$ 处具有相同的\omterm{def:b3:submanifolds:tangent}{切空间}：即反对称矩阵。

\textbf{8.} $A_u^2v = u\wedge(u\wedge v) = \langle u,
v\rangle u - v$（$u$ 为单位向量），故 $A_u^3v = u\wedge(\langle
u,v\rangle u - v) = -u\wedge v$：$A_u^3 = -A_u$。将
指数级数按幂次的剩余类分组，模关系
$A^3 = -A$：
\[
\eu^{\theta A_u} = I + \Bigl(\theta - \frac{\theta^3}{3!} +
\cdots\Bigr)A_u + \Bigl(\frac{\theta^2}{2!} -
\frac{\theta^4}{4!} + \cdots\Bigr)A_u^2
= I + \sin\theta\,A_u + (1 - \cos\theta)\,A_u^2 .
\]
对 $u$：$A_uu = 0$，故不动。对 $w \perp u$：$\eu^{\theta
A_u}w = w + \sin\theta\,u\wedge w + (1 - \cos\theta)(-w) =
\cos\theta\,w + \sin\theta\,u\wedge w$：即轴为
$u$、角为 $\theta$ 的旋转——罗德里格斯公式。每个旋转都具有
此形式（第 6 题）：由反对称
矩阵出发，$\exp$ 满射到 $SO(3)$。

\textbf{9.} 取 $q(t) = \cos\frac t2 + \sin\frac t2\,u$：
第 5 题表明 $\rho_{q(t)}$ 是轴为 $u$
且角为 $t$ 的旋转，即 $\rho_{q(t)} = \eu^{tA_u}$，其
导数在 $t = 0$ 处为 $A_u$。四元数以一半的
角度运动——这正是双覆盖的解析痕迹。

\textbf{10.} $\rho_{q(t)}$ 是绕 $\mathrm k$
转角 $t$ 的旋转：当 $t = 2\pi$ 时回到单位矩阵——
即 $SO(3)$ 中的一个回路。其提升满足 $q(2\pi) = \cos\pi =
-1 = -q(0)$：提升路径\emph{不}闭合；只有在
$t = 4\pi$ 时 $q$ 才回到 $1$。解释：由整圈旋转
构成的回路在 $SO(3)$ 中不可缩——其
提升终点位于覆盖的另一叶上——而
双圈回路则可缩；一个以绳带与周围相连的物体
（即皮带技巧）在
转过 $4\pi$ 后（而非 $2\pi$ 后）恢复无扭状态。旋转群记住
了整圈的奇偶性；$S^3$ \omterm{def:b3:conformal:simplyconnected}{单连通}，正是这种记忆
所在之处。

\textbf{11.} 四分之一转：$q_{\mathrm i} =
\cos\frac\pi4 + \sin\frac\pi4\,\mathrm i$, $q_{\mathrm j} =
\cos\frac\pi4 + \sin\frac\pi4\,\mathrm j$。乘积
（先作用 $\mathrm j$ 转动）：
\[
q_{\mathrm i}q_{\mathrm j} = \tfrac12(1 + \mathrm i)(1 +
\mathrm j) = \tfrac12\bigl(1 + \mathrm i + \mathrm j +
\mathrm k\bigr),
\]
模为 $1$，由 $\cos\frac\theta2 = \frac12$ 得 $\theta =
\frac{2\pi}3$，轴为 $u = \frac{\mathrm i + \mathrm j +
\mathrm k}{\sqrt3}$（纯虚部单位化）。两次依次进行的
绕相互正交轴的四分之一转合成一个
绕立方体主对角线的 $120^\circ$ 转动——运算
量仅相当于四次实数乘法，正交性
被精确保持：这正是飞行软件与图形引擎
借助四元数复合旋转的原因。

\textbf{12.} 由表可知 $\mathrm{ji} = -\mathrm k$ 且
$\mathrm{ki} = \mathrm j$，故
\[
q\,\mathrm i = a\mathrm i - b + c(\mathrm{ji}) +
d(\mathrm{ki}) = -b + a\mathrm i + d\mathrm j - c\mathrm k .
\]
乘以 $\bar q = a - b\mathrm i - c\mathrm j -
d\mathrm k$，并利用数量--向量法则 $(t_1 + p_1)(t_2 +
p_2) = t_1t_2 - \langle p_1, p_2\rangle + t_1p_2 + t_2p_1 +
p_1\wedge p_2$，其中 $p_1 = (a, d, -c)$ 且 $p_2 = (-b, -c,
-d)$：数量部分为 $-ab - (-ab) = 0$（纯虚，
必然如此），向量部分为
\[
b(b, c, d) + a(a, d, -c) + (-c^2 - d^2,\ bc + ad,\ bd - ac)
= \bigl(a^2 + b^2 - c^2 - d^2,\ 2(bc + ad),\ 2(bd -
ac)\bigr):
\]
即 $R_q$ 的第一列。循环映射 $\sigma(\mathrm i)
= \mathrm j$, $\sigma(\mathrm j) = \mathrm k$,
$\sigma(\mathrm k) = \mathrm i$ 保持关系
$\mathrm i^2 = \mathrm j^2 = \mathrm k^2 = \mathrm{ijk} =
-1$（词 $\mathrm{ijk}$ 在循环意义下不变，至多相差
关系 $\mathrm{ijk} = \mathrm{jki}$，而该关系在
任意环中成立：将 $\mathrm{ijk} = -1$ 用可逆元
$\mathrm i$ 作共轭即得），故 $\sigma$ 可延拓为 $\R$ 代数的
自同构，且 $\sigma(\rho_q(v)) =
\rho_{\sigma(q)}(\sigma(v))$。展开可知
$\mathrm j$ 的像即第一列公式经
代换 $(b, c, d) \to (c, d, b)$ 并将基
重新标记为 $\mathrm i \to \mathrm j \to \mathrm k \to \mathrm
i$，恰为所示的第二列；再作一次
即得第三列。

\textbf{13.} 对角线求和得 $\operatorname{tr}R_q =
3a^2 - (b^2 + c^2 + d^2) = 4a^2 - 1$（单位模长），并取
$a = \cos\frac\theta2$，则 $4\cos^2\frac\theta2 - 1 = 1 +
2\cos\theta$。反对称部分的三个独立
元素为 $R_q - R_q^{\mathsf T}$ 的 $4ab, 4ac, 4ad$（位于
位置 $(3,2), (1,3), (2,1)$），故 $\frac12(R_q -
R_q^{\mathsf T}) = A_m$，其中 $m = 2a\,(b, c, d) =
2\cos\frac\theta2\sin\frac\theta2\,u = \sin\theta\,u$。
算法：$\theta = \arccos\frac{\operatorname{tr}R -
1}{2} \in \intcc0\pi$；若 $0 < \theta < \pi$，则 $u$ 由
$\frac{R - R^{\mathsf T}}{2\sin\theta}$ 读出，并取 $q =
\pm(\cos\frac\theta2 + \sin\frac\theta2 u)$；若 $\theta =
0$，则 $q = \pm1$。对 $\theta = \pi$：$a = 0$，由罗德里格斯
公式（第 8 题）得 $R = I + 2A_u^2 = I + 2(uu^{\mathsf T} -
I) = 2uu^{\mathsf T} - I$，即 $R + I = 2uu^{\mathsf T}$；
$R + I$ 的任一非零列单位化后为 $\pm u$，且
$q = \pm u$。

\textbf{14.} 取 $a = b = c = d = \frac12$ 时，对角元
全为零，$2(bc - ad) = 0$，$2(bd + ac) = 1$，$2(bc +
ad) = 1$，$2(cd - ab) = 0$，$2(bd - ac) = 0$，$2(cd + ab) =
1$：
\[
R_q = \begin{pmatrix} 0 & 0 & 1\\ 1 & 0 & 0\\ 0 & 1 & 0
\end{pmatrix},
\]
即循环置换 $e_1 \to e_2 \to e_3 \to e_1$。其
迹为 $0 = 1 + 2\cos\theta$，故 $\theta = \frac{2\pi}3$，
且它固定 $(1,1,1)$：即绕
主对角线的 $120^\circ$ 旋转——恰为第 11 题的答案，现
可见它即循环坐标轴的矩阵。

\textbf{15.} 在基上验算可得 $\Phi(\bar q) =
\Phi(q)^*$（$\bar q$ 的矩阵满足 $\alpha' =
\bar\alpha$，$\beta' = -\beta$，恰是
$\bigl(\begin{smallmatrix}\alpha & \beta\\
-\bar\beta & \bar\alpha\end{smallmatrix}\bigr)$）的共轭转置。由此
$\Phi(q)^*\Phi(q) = \Phi(\bar qq) = N(q)I$，且 $\det\Phi(q)
= \abs\alpha^2 + \abs\beta^2 = N(q)$：对 $q \in S^3$，
$\Phi(q) \in SU(2)$，且 $\Phi$ 为单同态
（第 1 题）。满射性：设 $U =
\bigl(\begin{smallmatrix}\alpha & \beta\\ \gamma &
\delta\end{smallmatrix}\bigr)$，满足 $\det U = 1$；则
$U^{-1} = \bigl(\begin{smallmatrix}\delta & -\beta\\
-\gamma & \alpha\end{smallmatrix}\bigr)$，且 $U^{-1} = U^*
= \bigl(\begin{smallmatrix}\bar\alpha & \bar\gamma\\
\bar\beta & \bar\delta\end{smallmatrix}\bigr)$ 迫使
$\delta = \bar\alpha$，$\gamma = -\bar\beta$，进而 $1 =
\det U = \abs\alpha^2 + \abs\beta^2$：$U = \Phi(q)$，对应的
单位四元数 $q$ 的坐标为 $\alpha = a + \iu b$，
$\beta = c + \iu d$。于是 $S^3 \cong SU(2)$。

\textbf{16.} 实数是中心元，又由 $N(p) = 1$ 得
$\overline{pq\bar p} = p\bar q\bar p$，故 $pq\bar p +
\overline{pq\bar p} = p(q + \bar q)\bar p = q + \bar q$：
实部是不变量。反之设
$\operatorname{Re}q = \operatorname{Re}q' = a$；
则纯部的模同为 $\sqrt{1 - a^2} = s$。若 $s =
0$，则 $q = q' = \pm1$。若 $s > 0$，记 $q = a + su$，$q' =
a + su'$，其中 $u, u'$ 为单位纯四元数；由第 6 题，
存在把 $u'$ 变为 $u$ 的旋转，即 $p \in S^3$，使
$\rho_p(u') = u$，于是 $pq'\bar p = a + s\rho_p(u') =
q$。因此 $S^3$ 的\omterm{ex:b3:groups:actions}{共轭类}为 $\{1\}$，$\{-1\}$，
以及对每个 $a \in \intoo{-1}1$，$2$ 维球面 $\{a + su :
u \text{ 为单位纯}\}$，半径为 $s$。在 $\Phi$ 下，
$\operatorname{tr}\Phi(q) = \alpha + \bar\alpha =
2\operatorname{Re}q$：$SU(2)$ 的\omterm{ex:b3:groups:actions}{共轭类}恰为迹的
水平集。就投影而言，若 $q' = pq\bar p$，则
$\rho_{q'} = \rho_p\rho_q\rho_p^{-1}$；反之
$\rho_{q'} = \rho_p\rho_q\rho_p^{-1} = \rho_{pq\bar p}$
迫使 $q' = \pm pq\bar p$（由核的性质），故
$\operatorname{Re}q' = \pm\operatorname{Re}q$，即
$\cos\frac{\theta'}2 = \abs{\operatorname{Re}q'} =
\abs{\operatorname{Re}q} = \cos\frac\theta2$，对
$\intcc0\pi$ 中的角，共轭旋转的角相等。
反之，角相等时，可取代表元，使其具有
相同的非负实部，由上可知二者共轭：在
$SO(3)$ 中，旋转的\omterm{ex:b3:groups:actions}{共轭类}恰由其
角决定。

\textbf{17.} $N$ 是乘性的，故 $\abs\cdot$ 是
$\mathbb H \cong \R^4$ 上的乘性范数，且
$\abs{q^k} = \abs q^k$：级数 $\sum q^k/k!$ 在
有限维（从而完备）空间中绝对收敛，
被 $\sum\abs q^k/k! = \eu^{\abs q}$ 所控制。对
单位纯四元数 $u$：$u^2 = -1$，故 $(\theta u)^{2m} =
(-1)^m\theta^{2m}$，且 $(\theta u)^{2m+1} =
(-1)^m\theta^{2m+1}u$；拆分级数，
\[
\exp(\theta u) = \sum_m\frac{(-1)^m\theta^{2m}}{(2m)!} +
u\sum_m\frac{(-1)^m\theta^{2m+1}}{(2m+1)!} = \cos\theta +
\sin\theta\,u .
\]
任一 $q \in S^3$ 都可写成 $\cos\alpha + \sin\alpha\,u$，其中
$\alpha \in \intcc0\pi$（第 5 题）：$q = \exp(\alpha u)$，
故 $\exp(P) = S^3$。最后 $\exp(su) = \cos s + \sin s\,u
= q(2s)$，依第 9 题的记号，且 $\rho_{q(t)} =
\eu^{tA_u}$ 在此，即 $\rho_{\exp(su)} = \eu^{2sA_u}$。

\textbf{18.} 利用乘法表逐分量展开 $vw$：
乘积 $\mathrm i\cdot\mathrm i = -1$，\dots\ 给出
数量 $-(v_1w_1 + v_2w_2 + v_3w_3)$，而混合
乘积（$\mathrm{ij} = \mathrm k$，$\mathrm{ji} =
-\mathrm k$，\dots）给出向量 $(v_2w_3 - v_3w_2,\
v_3w_1 - v_1w_3,\ v_1w_2 - v_2w_1)$：$vw = -\langle v,
w\rangle + v\wedge w$。减去反序乘积：
$vw - wv = 2\,v\wedge w$（数量部分相消，
向量积相加）。对矩阵，利用 $a\wedge(b\wedge
c) = b\langle a, c\rangle - c\langle a, b\rangle$：
\[
[A_v, A_w]x = v\wedge(w\wedge x) - w\wedge(v\wedge x)
= w\langle v, x\rangle - v\langle w, x\rangle
= (v\wedge w)\wedge x = A_{v\wedge w}x .
\]
求导：$\overline{\exp(tv)} = \exp(-tv)$（共轭
连续，且使纯四元数变号），故
\[
\frac{\dd}{\dd t}\Bigr|_{t=0}\exp(tv)\,x\,\exp(-tv) = vx -
xv = 2\,v\wedge x = 2A_vx :
\]
微分为 $v \mapsto 2A_v$，是一个线性双射，
从 $P$ 满射到反对称矩阵上，且 $[2A_v, 2A_w]
= 4A_{v\wedge w} = 2A_{2v\wedge w} = 2A_{[v,w]}$ 表明它
将四元数的\omterm{def:b3:groups:derived}{换位子}映为矩阵的
\omterm{def:b3:groups:derived}{换位子}。

\textbf{19.} 作用 $\rho$ 得：$\rho(\varepsilon(t)) =
\rho(s(R(t)))\,\rho(q(t))^{-1} = R(t)R(t)^{-1} =
\operatorname{id}$，故 $\varepsilon(t) \in \ker\rho =
\{\pm1\}$（第 4 题）。
由\omterm{def:b3:topology:continuity}{连续映射} $t \mapsto s(R(t))$ 与 $t \mapsto q(t)^{-1} =
\bar{q(t)}$ 之积，$\varepsilon$ 在连通的
区间 $\intcc0{2\pi}$ 上连续，取值于离散的
二元集 $\{\pm1\}$：它为常值，记作 $\varepsilon(t) \equiv
\varepsilon$。但 $R(0) = R(2\pi) = I$，故 $s(R(0)) =
s(R(2\pi))$，而 $s(R(0)) = \varepsilon\,q(0) =
\varepsilon$，且 $s(R(2\pi)) = \varepsilon\,q(2\pi) =
-\varepsilon$：矛盾。不存在连续
截面：符号的不定性 $\pm q$ 是整体的，而非
某一特定公式的缺陷。

\textbf{20.} \emph{满射：}每个 $R \in SO(3)$ 都可写成
$\eu^{\theta A_u}$ 的形式，其中 $u$ 为某单位向量，$\theta \in
\intcc0{2\pi}$（第 6 题与第 8 题）；若 $\theta > \pi$，
由罗德里格斯公式得 $\eu^{\theta A_u} = \eu^{(2\pi -
\theta)A_{-u}}$（两者均等于 $I + \sin\theta A_u + (1 -
\cos\theta)A_u^2$，且 $A_{-u} = -A_u$，其中 $\sin(2\pi -
\theta) = -\sin\theta$，$\cos(2\pi - \theta) =
\cos\theta$），于是 $R = E(v)$，其中 $\norm v \leq \pi$。
\emph{\omterm{def:b3:topology:interior}{内部}单射：}若 $E(v) = E(v') \neq I$，且
$\norm v, \norm{v'} < \pi$，由第 13 题得同一角度
$\theta = \norm v = \norm{v'} \in \intoo0\pi$，由迹确定，
并因 $\sin\theta \neq 0$，同一轴由
反对称部分确定：$v = v'$；又 $E(v) = I$
在开球上蕴含 $\theta \in \{0\}$。\emph{边界：}
$E(\pi u) = I + 2A_u^2 = 2uu^{\mathsf T} - I$ 依赖于
$u$ 且仅经由 $uu^{\mathsf T}$，故 $E(\pi u) =
E(-\pi u)$；反之，$2uu^{\mathsf T} - I =
2u'u'^{\mathsf T} - I$ 作用于 $u$ 得 $u = \langle
u', u\rangle u'$，故 $u' = \pm u$。\omterm{def:b3:topology:interior}{内部}与边界
互不相交（迹 $> -1$，对比 $= -1$）。于是 $E$ 诱导一个
从带对径粘合之闭球
——紧——到 $SO(3)$ 上的连续双射：它是\omterm{def:b3:topology:continuity}{同胚}，且 $SO(3)
\cong \mathbb{RP}^3$。从 $\pi u$ 到 $-\pi u$ 的直径
两端已粘合：它是 $SO(3)$ 中的一个圈，其
$E$ 描述与第 10 题所述绕轴旋转之族相吻合，
即绕 $u$ 扫过整整一圈的旋转。

\textbf{21.} $\rho$ 为同态，且 $\rho_p^{-1} =
\rho_{p^{-1}} = \rho_{\bar p}$，故 $\rho_p\rho_q\rho_p^{-1}
= \rho_{pq\bar p}$；由第 16 题，共轭一个旋转
保持其角度，并将其轴经 $\rho_p$ 旋转。
对合：$\rho_q^2 = \rho_{q^2} = \operatorname{id}$
当且仅当 $q^2 = \pm1$。若 $q^2 = 1$，则 $(q - 1)(q + 1) = q^2 -
1 = 0$（标量位于中心，故该分解有效），且
$q = \pm1$（在除环 $\mathbb H$ 中），于是 $\rho_q
= I$，应予排除；当 $q^2 = -1$ 时，取 $q = a +
su$ 得 $a^2 - s^2 + 2as\,u = -1$，故 $a = 0$，$s = 1$：
$q$ 为单位纯四元数 $w$，而 $\rho_w$ 为绕
$w$ 的半转（角度 $\pi$，第 5 题）。中心：若 $\rho_q$
与每个 $\rho_p$ 交换，则 $\rho_{pq\bar p} = \rho_q$，故
$pq\bar p = \varepsilon(p)\,q$，其中 $\varepsilon(p) \in
\{\pm1\}$；映射 $p \mapsto \varepsilon(p) = (pq\bar p)q^{-1}$
在连通的 $S^3$ 上连续，且取值 $1$，即在 $p =
1$ 时，故 $\equiv 1$：$q$ 与 $S^3$ 中所有元素可交换，从而
与整个 $\mathbb H$ 可交换（经数乘放缩），故 $q \in \R \cap S^3 =
\{\pm1\}$（第 1 题），且 $\rho_q = I$：中心
是平凡的。

\textbf{22.} 由于 $u \perp w$ 为单位纯四元数，$uw =
u\wedge w$ 为纯四元数（第 18 题），故
\[
w' = qw = \cos\tfrac\theta2\,w + \sin\tfrac\theta2\,u\wedge
w
\]
它为纯四元数，范数为 $\abs q\abs w = 1$。于是 $w'w = qw\cdot w
= qw^2 = -q$，且
\[
\rho_{w'}\rho_w = \rho_{w'w} = \rho_{-q} = \rho_q .
\]
两轴 $w$ 与 $w' = \cos\frac\theta2 w +
\sin\frac\theta2(u\wedge w)$ 均位于与旋转轴正交的平面 $u^\perp$
内，且 $\langle w', w\rangle
= \cos\frac\theta2$：二者所成的角为
$\frac\theta2$。这即经典的生成：两个
绕交角为 $\frac\theta2$ 之轴的半转
复合为角度为 $\theta$ 的旋转，其轴为二者的
公垂线。

\textbf{23.} 紧性即第 7 题。\omterm{def:b3:topology:pathconnected}{道路}连通性：
$SO(3) = \rho(S^3)$ 为
\omterm{def:b3:topology:pathconnected}{道路}连通球面的连续像；另外，对 $R = \eu^{A}$
其中 $A$ 为反对称（第 8 题），$t \mapsto \eu^{tA}$
是 $SO(3)$ 中从 $I$ 到 $R$ 的\omterm{def:b3:topology:pathconnected}{道路}（正交，因
$(\eu^{tA})^{\mathsf T} = \eu^{-tA}$，行列式 $1$ 由
从 $t = 0$ 出发的连续性得出）。对 $O(3)$：$\det$ 连续
满射到 $\{\pm1\}$，故 $O(3)$ 不连通，$O(3) = SO(3)
\sqcup D\,SO(3)$，其中 $D$ 为任一固定的、满足 $\det D = -1$
的元素（例如 $D = -I$），而左乘 $D$ 是一个
\omterm{def:b3:topology:continuity}{同胚}：恰有两个\omterm{def:b3:topology:components}{连通分支}，每个都是
$SO(3)$ 的副本。二对一的 $\rho$，由
第 19 题知其无截面，因而是一个连通紧
群的真正双覆盖，覆盖空间为\omterm{def:b3:conformal:simplyconnected}{单连通}的 $S^3$——这正是
皮带戏法背后的几何。

\textbf{24.} (i) 矩阵：
\[
R_1 = \begin{pmatrix} 0 & -1 & 0\\ 1 & 0 & 0\\ 0 & 0 & 1
\end{pmatrix},
\quad
R_2 = \begin{pmatrix} 1 & 0 & 0\\ 0 & 0 & -1\\ 0 & 1 & 0
\end{pmatrix},
\quad
R_2R_1 = \begin{pmatrix} 0 & -1 & 0\\ 0 & 0 & -1\\ 1 & 0 &
0\end{pmatrix}.
\]
由迹 $0 = 1 + 2\cos\theta$ 得 $\cos\theta = -\frac12$：
$\theta = \frac{2\pi}3$。反对称部分 $\frac{R -
R^{\mathsf T}}2$ 的元素按 $\sin\theta\,(v_3,
-v_2, v_1)$ 的方式给出轴：此处 $\frac{R - R^{\mathsf T}}2 =
\frac12\bigl(\begin{smallmatrix}0 & -1 & -1\\ 1 & 0 & -1\\ 1
& 1 & 0\end{smallmatrix}\bigr)$，据（第五部分的
字典 $A_v$），可读出 $v\sin\theta = \frac12(1, -1, 1)$；
由 $\sin\frac{2\pi}3 = \frac{\sqrt3}2$，得 $v = \frac1{\sqrt3}(1,
-1, 1)$。
(ii) 四元数：$q_1 = \cos\frac\pi4 + \sin\frac\pi4\,k =
\frac{\sqrt2}2(1 + k)$，$q_2 = \frac{\sqrt2}2(1 + i)$，且
\[
q_2q_1 = \tfrac12(1 + i)(1 + k) = \tfrac12(1 + k + i + ik)
= \tfrac12\bigl(1 + i - j + k\bigr)
\]
（$ik = -j$）。于是 $\cos\frac\theta2 = \frac12$：$\theta =
\frac{2\pi}3$，向量部分 $\frac12(i - j + k)$ 的
方向为 $\frac1{\sqrt3}(1, -1, 1)$——答案相同，
四元数路线仅需一行乘法，
而非矩阵乘积：这正是
飞行软件在 $S^3$ 中合成姿态的实际原因。

\textbf{25.} $I + K$ 可逆：由 $(I + K)v = 0$ 得 $0 =
\langle v, v\rangle + \langle v, Kv\rangle = \norm v^2$
（反对称性使第二项为零）：$v = 0$。
$C = C(K)$ 的正交性：利用 $(I \pm K)^{\mathsf T} = I \mp K$ 及
事实：四个矩阵 $I \pm K$、$(I \pm K)^{-1}$
两两可交换（为 $K$ 的多项式，并取极限）：
\[
C^{\mathsf T}C = (I + K)^{-\mathsf T}(I - K)^{\mathsf T}
(I - K)(I + K)^{-1}
= (I - K)^{-1}(I + K)(I - K)(I + K)^{-1} = I .
\]
行列式：$\det(I - K) = \det\bigl((I - K)^{\mathsf
T}\bigr) = \det(I + K)$，故 $\det C = 1$：$C \in SO(n)$。
$-1$ 不是特征值：$Cv = -v$ 意味着 $(I - K)w = -(I + K)w$
令 $w = (I + K)^{-1}v$，即 $2w = 0$：$v = 0$。
求逆：由 $C(I + K) = I - K$，解得 $K(I + C) = I -
C$；由于 $-1 \notin \operatorname{Sp}C$，$I + C$
可逆，且 $K = (I - C)(I + C)^{-1}$，它在
$C$ 正交且无特征值
$-1$ 时反对称（转置此式，并用 $C^{\mathsf T} =
C^{-1}$：$K^{\mathsf T} = (I - C^{-1})(I + C^{-1})^{-1} =
(C - I)(C + I)^{-1} = -K$）。两映射由构造而
互逆：这是一个全局有理参数化，
覆盖 $SO(n)$ 中避开特征值 $-1$ 的稠密开部分
——无需级数，也无需三角函数；在维数 $3$ 时，
它实为半角代换 $K = \tan\frac\theta2\,A_v$ 的
伪装。
\end{solution}
